% BHU PYQ -- Professional Exam Template
\documentclass[11pt,a4paper]{article}

\usepackage[a4paper,top=2.5cm,bottom=2.5cm,left=2.8cm,right=2.8cm,headheight=14pt]{geometry}
\usepackage{amsmath,amssymb}
\usepackage[T1]{fontenc}
\usepackage[utf8]{inputenc}
\usepackage{lmodern}
\usepackage{microtype}
\usepackage{fancyhdr}
\usepackage{enumitem}
\usepackage{listings}
\usepackage{hyperref}
\usepackage{lastpage}
\usepackage{parskip}

\hypersetup{colorlinks=true,urlcolor=black,linkcolor=black,
  pdftitle={BCS-201: Digital Logic Circuits -- BHU PYQ 2013-14}}

\pagestyle{fancy}
\fancyhf{}
\renewcommand{\headrulewidth}{0.4pt}
\renewcommand{\footrulewidth}{0.4pt}
\fancyhead[L]{\small\textit{Banaras Hindu University}}
\fancyhead[R]{\small\textit{BCS-201: Digital Logic Circuits}}
\fancyfoot[C]{\small Page \thepage\ of \pageref{LastPage}}

\lstset{
  basicstyle=\small\ttfamily,
  frame=single,
  breaklines=true,
  columns=flexible,
  keepspaces=true
}

\newlist{parts}{enumerate}{2}
\setlist[parts,1]{label=(\alph*),leftmargin=*,itemsep=4pt,topsep=3pt}
\setlist[parts,2]{label=(\roman*),leftmargin=*,itemsep=2pt,topsep=2pt}
\newcommand{\pts}[1]{\hfill\textit{[#1]}}

\begin{document}
\begin{center}
  {\large\bfseries BANARAS HINDU UNIVERSITY}\\[2pt]
  {\normalsize B.Sc.\ (Hons.) Semester II Examination, 2013-14}\\[6pt]
  \rule{0.8\linewidth}{0.5pt}\\[6pt]
  {\large\bfseries Paper: BCS-201 --- Digital Logic Circuits}\\[4pt]
  {\normalsize Time: Three hours \hfill Maximum Marks: 70}
\end{center}

\rule{\linewidth}{0.4pt}

\smallskip
\noindent\textit{Instructions: Answer any \textbf{five} questions, including Question~1, which is compulsory.}

\bigskip

%% ─── Question 1 (Compulsory) ──────────────────────────────────────────────────
\noindent\textbf{Question 1.}
\begin{parts}
  \item Perform the following conversions. Show the detailed calculation being done for conversions:\pts{3$\times$2 = 6}
  \begin{parts}
    \item $(673.23)_{10} = (\quad)_2$
    \item $(1101101.111)_2 = (\quad)_{10}$
    \item $(2AC5.D)_{16} = (\quad)_8$
  \end{parts}
  \item Represent the decimal number $3620$ in BCD and excess-3 code.\pts{2}
  \item What is race around condition?\pts{2}
  \item List the names of any four digital logic families.\pts{2}
  \item What is the advantage of PLA over ROM?\pts{2}
\end{parts}

\bigskip

%% ─── Question 2 ───────────────────────────────────────────────────────────────
\noindent\textbf{Question 2.}
\begin{parts}
  \item Perform the following binary operations:\pts{3}
  \begin{parts}
    \item $111 \times 1101$
    \item $11011001 \div 1100$
  \end{parts}
  \item Subtract $3$ from $-12$ using 2's complement method.\pts{2}
  \item Perform the addition of the following numbers in the given base without converting to decimal:\pts{4}
  \begin{parts}
    \item $(1230)_4 + (23)_4$
    \item $(296)_{12} + (57)_{12}$
  \end{parts}
  \item What do you understand by floating point representation? Represent $(4.21875)_{10}$ in binary floating point notation.\pts{5}
\end{parts}

\bigskip

%% ─── Question 3 ───────────────────────────────────────────────────────────────
\noindent\textbf{Question 3.}
\begin{parts}
  \item Simplify the following Boolean expressions to a minimum number of literals by using Boolean algebra theorems:\pts{4}
  \begin{parts}
    \item $xyz - x'y - xyz'$
    \item $BC + AC' + AB + BCD$
  \end{parts}
  \item What is the difference between canonical form and standard form? Which form is preferable when implementing a Boolean function with gates? Which form is obtained when reading a function from a truth table?\pts{5}
  \item Express the following function in a sum of minterms and product of maxterms:
  \[ F(w, x, y, z) = y'z + wxy' + wxz' + w'x'z \]\pts{5}
\end{parts}

\bigskip

%% ─── Question 4 ───────────────────────────────────────────────────────────────
\noindent\textbf{Question 4.}
\begin{parts}
  \item What are the differences between combinational circuit and sequential circuit? Explain with their block diagrams. List the names of any two combinational and any two sequential circuits.\pts{4}
  \item Implement the following function with NAND gates (with no more than six gates, each having three inputs). Assume that both the normal and complement inputs are available.
  \[ F = BD + BCD + AB'C'D' + A'B'CD' \]\pts{5}
  \item Simplify the Boolean function $F$ using the don't care condition $d$, in its minimum sum of product form:
  \begin{align*}
    F(w, x, y, z) &= w'x'y'z' + w'x'yz' + w'xy'z + w'xyz' \\
    d(w, x, y, z) &= w'x'y'z + yz' + wyz
  \end{align*}\pts{5}
\end{parts}

\bigskip

%% ─── Question 5 ───────────────────────────────────────────────────────────────
\noindent\textbf{Question 5.}
\begin{parts}
  \item Design a combinational circuit that accepts a three bit number and generates an output binary number equal to the square of the input number.\pts{5}
  \item What is a decoder? Design a full adder circuit with a decoder and two OR gates.\pts{4}
  \item What is meant by a priority encoder? How does a priority encoder differ from an ordinary encoder? Explain with suitable example.\pts{5}
\end{parts}

\bigskip

%% ─── Question 6 ───────────────────────────────────────────────────────────────
\noindent\textbf{Question 6.}
\begin{parts}
  \item Define Multiplexer. Implement an 8x1 multiplexer using two 4x1 multiplexers and a 2x1 multiplexer.\pts{5}
  \item Implement a 3-bit binary to gray code converter using PLA as follows:
  \begin{parts}
    \item Make a truth table.\pts{2}
    \item Simplify using K-map.\pts{2}
    \item Design circuit using PLA; and\pts{3}
    \item Specify the size of PLA.\pts{2}
  \end{parts}
\end{parts}

\bigskip

%% ─── Question 7 ───────────────────────────────────────────────────────────────
\noindent\textbf{Question 7.}
\begin{parts}
  \item Implement JK flip flop using D-flip flop.\pts{4}
  \item Define irregular counter. Design a counter with the following binary sequence: $0, 1, 3, 7, 6, 4$ and repeat. Use T flip-flops.\pts{5}
  \item Define register. By using registers, design a sequential circuit whose state table is listed as follows:\pts{5}
  
  \begin{center}
  \begin{tabular}{|cc|c|cc|c|}
  \hline
  \multicolumn{2}{|c|}{\textbf{Present State}} & \textbf{Input} & \multicolumn{2}{c|}{\textbf{Next State}} & \textbf{Output} \\
  \hline
  $A_1$ & $A_2$ & $x$ & $A_1$ & $A_2$ & $y$ \\
  \hline
  0 & 0 & 0 & 0 & 0 & 0 \\
  0 & 0 & 1 & 0 & 1 & 0 \\
  0 & 1 & 0 & 0 & 1 & 0 \\
  0 & 1 & 1 & 0 & 0 & 1 \\
  1 & 0 & 0 & 1 & 0 & 0 \\
  1 & 0 & 1 & 0 & 1 & 0 \\
  1 & 1 & 0 & 1 & 1 & 0 \\
  1 & 1 & 1 & 0 & 0 & 1 \\
  \hline
  \end{tabular}
  \end{center}
\end{parts}

\vfill
\rule{\linewidth}{0.4pt}
\begin{center}\small
  \href{https://bhu-exam.vercel.app/}{Click here to visit our website}\\[4pt]
  \href{https://me-aryan.vercel.app/}{Click here to visit our contributor}\\[4pt]
  \href{https://quashan-me.vercel.app/}{Click here to visit our contributor}
\end{center}
\end{document}